\chapter{System Design and Methodology}

\section{System Overview}

ShadeAI is designed as a modular, pipeline-based system that processes facial photographs through six distinct stages to produce foundation shade recommendations. The system architecture emphasizes:

\begin{itemize}
    \item \textbf{Modularity:} Each component is independent and can be tested, modified, or replaced without affecting others
    \item \textbf{Scalability:} The system can process single images or batches efficiently
    \item \textbf{Accessibility:} No training data or GPU hardware required
    \item \textbf{Accuracy:} Research-grade performance using classical computer vision and color science
    \item \textbf{Usability:} Multiple interfaces for different use cases
\end{itemize}

\begin{figure}[h]
\centering
\begin{verbatim}
Input Image
    ↓
[1] Preprocessing (White Balance, Color Spaces)
    ↓
[2] Face Detection (Haar Cascade)
    ↓
[3] Skin Segmentation (HSV + YCbCr)
    ↓
[4] Feature Extraction (ITA°, LAB Statistics)
    ↓
[5] Classification (MST, Undertone)
    ↓
[6] Shade Recommendation (CIEDE2000)
    ↓
Output (Top 5 Shades)
\end{verbatim}
\caption{ShadeAI Processing Pipeline}
\end{figure}

\section{System Architecture}

\subsection{High-Level Architecture}

The system follows a three-tier architecture:

\begin{enumerate}
    \item \textbf{Presentation Layer:}
    \begin{itemize}
        \item Command-line interface (shadeai\_cli.py)
        \item Web application (streamlit\_app.py)
        \item REST API (backend/api/routes.py)
    \end{itemize}
    
    \item \textbf{Business Logic Layer:}
    \begin{itemize}
        \item Pipeline modules (backend/pipeline/)
        \item Utility functions (backend/utils/)
        \item Photo quality checker
    \end{itemize}
    
    \item \textbf{Data Layer:}
    \begin{itemize}
        \item Foundation shade database (JSON)
        \item Pre-trained models (Haar Cascade)
        \item Configuration files
    \end{itemize}
\end{enumerate}

\subsection{Module Organization}

\begin{verbatim}
shadeai/
├── backend/
│   ├── pipeline/
│   │   ├── preprocessor.py          # Module 1
│   │   ├── face_detector.py         # Module 2
│   │   ├── skin_segmentor.py        # Module 3
│   │   ├── feature_extractor.py     # Module 4
│   │   ├── classifier.py            # Module 5
│   │   └── shade_recommender.py     # Module 6
│   ├── utils/
│   │   ├── ciede2000.py
│   │   ├── color_utils.py
│   │   └── photo_quality.py
│   ├── api/
│   │   └── routes.py
│   ├── data/
│   │   └── foundation_shades.json
│   └── main.py
├── shadeai_cli.py
├── streamlit_app.py
└── requirements.txt
\end{verbatim}

\section{Pipeline Modules}

\subsection{Module 1: Image Preprocessing}

\subsubsection{Purpose}
The preprocessing module prepares input images for analysis by correcting color casts, converting color spaces, and validating image quality.

\subsubsection{White Balance Correction}

White balance correction removes color casts caused by different lighting conditions. ShadeAI implements the Gray World algorithm:

\textbf{Algorithm:}
\begin{enumerate}
    \item Calculate average RGB values: $\bar{R}, \bar{G}, \bar{B}$
    \item Compute gray value: $gray = (\bar{R} + \bar{G} + \bar{B}) / 3$
    \item Calculate scaling factors:
    \begin{align}
        k_R &= gray / \bar{R} \\
        k_G &= gray / \bar{G} \\
        k_B &= gray / \bar{B}
    \end{align}
    \item Apply corrections: $R' = R \times k_R$, $G' = G \times k_G$, $B' = B \times k_B$
    \item Clip values to [0, 255]
\end{enumerate}

\textbf{Rationale:} The Gray World algorithm assumes that the average color in a scene should be neutral gray. While simple, it is effective for most photographs and requires no training or calibration.

\subsubsection{Color Space Conversions}

The module converts images to multiple color spaces for different purposes:

\begin{itemize}
    \item \textbf{RGB:} Input format, used for display
    \item \textbf{CIELAB:} Used for ITA° calculation and color matching
    \item \textbf{HSV:} Used for skin segmentation
    \item \textbf{YCbCr:} Used for skin segmentation
\end{itemize}

\textbf{RGB to CIELAB Conversion:}
\begin{enumerate}
    \item Convert RGB to XYZ using standard transformation matrix
    \item Normalize XYZ by reference white point (D65)
    \item Apply nonlinear transformation:
    \begin{equation}
        f(t) = \begin{cases}
            t^{1/3} & \text{if } t > (6/29)^3 \\
            \frac{1}{3}(29/6)^2 t + 4/29 & \text{otherwise}
        \end{cases}
    \end{equation}
    \item Calculate LAB values:
    \begin{align}
        L^* &= 116 \cdot f(Y/Y_n) - 16 \\
        a^* &= 500 \cdot (f(X/X_n) - f(Y/Y_n)) \\
        b^* &= 200 \cdot (f(Y/Y_n) - f(Z/Z_n))
    \end{align}
\end{enumerate}

\subsection{Module 2: Face Detection}

\subsubsection{Purpose}
Face detection locates the face region in the image, which is essential for focusing analysis on facial skin rather than background or clothing.

\subsubsection{Haar Cascade Classifier}

ShadeAI uses OpenCV's pre-trained Haar Cascade classifier for face detection.

\textbf{Algorithm:}
\begin{enumerate}
    \item Convert image to grayscale
    \item Apply Haar Cascade classifier with parameters:
    \begin{itemize}
        \item scaleFactor = 1.1 (image pyramid scale)
        \item minNeighbors = 5 (minimum detections for confirmation)
        \item minSize = (30, 30) (minimum face size)
    \end{itemize}
    \item Select largest detected face (if multiple)
    \item Extract face region with 10\% padding
    \item Validate detection (minimum size, aspect ratio)
\end{enumerate}

\textbf{Advantages:}
\begin{itemize}
    \item Fast (real-time on CPU)
    \item No training required
    \item Robust to lighting variations
    \item Available in OpenCV
\end{itemize}

\textbf{Limitations:}
\begin{itemize}
    \item Less accurate for non-frontal faces
    \item May miss faces with occlusions
    \item Lower accuracy than deep learning methods
\end{itemize}

\subsubsection{Fallback Mechanisms}

If face detection fails:
\begin{enumerate}
    \item Try with relaxed parameters (minNeighbors = 3)
    \item Use entire image if still no detection
    \item Provide user feedback about photo quality
\end{enumerate}

\subsection{Module 3: Skin Segmentation}

\subsubsection{Purpose}
Skin segmentation identifies skin pixels within the face region, excluding eyes, eyebrows, mouth, hair, and background.

\subsubsection{Rule-Based Segmentation}

ShadeAI uses a combined HSV and YCbCr approach:

\textbf{HSV Thresholds:}
\begin{align}
    0° &\leq H \leq 50° \\
    0.2 &\leq S \leq 0.8 \\
    0.3 &\leq V \leq 1.0
\end{align}

\textbf{YCbCr Thresholds:}
\begin{align}
    77 &\leq Cb \leq 127 \\
    133 &\leq Cr \leq 173
\end{align}

\textbf{Algorithm:}
\begin{enumerate}
    \item Convert face region to HSV and YCbCr
    \item Apply thresholds to create binary masks
    \item Combine masks using logical AND
    \item Apply morphological operations:
    \begin{itemize}
        \item Opening (erosion + dilation) to remove noise
        \item Dilation to fill small gaps
    \end{itemize}
    \item Extract skin pixels using mask
    \item Validate (minimum 1000 pixels)
\end{enumerate}

\textbf{Rationale:} Combining HSV and YCbCr provides robustness to lighting variations while maintaining accuracy. HSV separates chromaticity from intensity, while YCbCr is optimized for skin detection based on empirical research.

\subsubsection{Morphological Operations}

Morphological operations clean up the segmentation mask:

\begin{itemize}
    \item \textbf{Opening:} Removes small noise regions
    \item \textbf{Dilation:} Fills small gaps and connects nearby regions
    \item \textbf{Kernel Size:} 5×5 pixels (empirically determined)
\end{itemize}

\subsection{Module 4: Feature Extraction}

\subsubsection{Purpose}
Feature extraction computes quantitative measures of skin color that can be used for classification and matching.

\subsubsection{ITA° Calculation}

The Individual Typology Angle (ITA°) is calculated from CIELAB values:

\begin{equation}
ITA° = \arctan\left(\frac{L^* - 50}{b^*}\right) \times \frac{180}{\pi}
\end{equation}

where:
\begin{itemize}
    \item $L^*$ is the mean lightness of skin pixels
    \item $b^*$ is the mean yellow-blue component
\end{itemize}

\textbf{Interpretation:}
\begin{itemize}
    \item Higher ITA° → Lighter skin tone
    \item Lower ITA° → Darker skin tone
    \item ITA° > 55° → Very fair (MST-01)
    \item ITA° < -30° → Deep (MST-10)
\end{itemize}

\subsubsection{CIELAB Statistics}

The module computes statistics for all three CIELAB components:

\begin{align}
    \bar{L^*} &= \frac{1}{N} \sum_{i=1}^{N} L_i^* \\
    \bar{a^*} &= \frac{1}{N} \sum_{i=1}^{N} a_i^* \\
    \bar{b^*} &= \frac{1}{N} \sum_{i=1}^{N} b_i^*
\end{align}

where $N$ is the number of skin pixels.

These statistics are used for:
\begin{itemize}
    \item ITA° calculation (tone)
    \item Undertone detection
    \item Color matching with foundation shades
\end{itemize}

\subsection{Module 5: Classification}

\subsubsection{MST Classification}

Skin tone is classified into one of 10 MST classes using ITA° thresholds:

\begin{table}[h]
\centering
\begin{tabular}{|c|l|c|}
\hline
\textbf{MST Class} & \textbf{Label} & \textbf{ITA° Range} \\
\hline
MST-01 & Porcelain & > 55° \\
MST-02 & Fair & 41° to 55° \\
MST-03 & Light & 28° to 41° \\
MST-04 & Light-Medium & 19° to 28° \\
MST-05 & Medium & 10° to 19° \\
MST-06 & Medium-Dark & 0° to 10° \\
MST-07 & Tan & -10° to 0° \\
MST-08 & Brown & -20° to -10° \\
MST-09 & Deep Brown & -30° to -20° \\
MST-10 & Deep & < -30° \\
\hline
\end{tabular}
\caption{ITA° Thresholds for MST Classification}
\end{table}

\textbf{Algorithm:}
\begin{verbatim}
if ITA > 55:
    class = MST-01
elif ITA > 41:
    class = MST-02
elif ITA > 28:
    class = MST-03
...
else:
    class = MST-10
\end{verbatim}

\subsubsection{Undertone Detection}

Undertone is detected using heuristic rules based on CIELAB $a^*$ and $b^*$ values:

\textbf{Algorithm:}
\begin{enumerate}
    \item Initialize score = 0
    \item If $b^* > 15$: score += 0.4 (yellow → warm)
    \item If $b^* < 10$: score -= 0.4 (blue → cool)
    \item If $a^* > 12$: score += 0.3 (red → warm)
    \item If $a^* < 8$: score -= 0.3 (pink → cool)
    \item If ITA° > 28: score += 0.2 (lighter → warm)
    \item If ITA° < 10: score -= 0.2 (darker → cool)
    \item Clip score to [-1.0, 1.0]
    \item Classify:
    \begin{itemize}
        \item score > 0.3 → Warm
        \item score < -0.3 → Cool
        \item otherwise → Neutral
    \end{itemize}
\end{enumerate}

\textbf{Rationale:} These heuristics are based on empirical observations and color theory. Warm undertones have higher yellow ($b^*$) and red ($a^*$) components, while cool undertones have lower values.

\subsection{Module 6: Shade Recommendation}

\subsubsection{Purpose}
The shade recommendation module matches the detected skin tone to foundation shades in the database using CIEDE2000 color difference.

\subsubsection{CIEDE2000 Implementation}

CIEDE2000 is the industry-standard metric for color difference. The implementation follows the CIE specification:

\textbf{Algorithm Steps:}
\begin{enumerate}
    \item Calculate chroma: $C = \sqrt{a^{*2} + b^{*2}}$
    \item Compute $\bar{C}$ and apply G correction to $a^*$
    \item Calculate hue angle: $h' = \arctan(b^*/a'^*)$
    \item Compute differences: $\Delta L'$, $\Delta C'$, $\Delta H'$
    \item Calculate weighting functions: $S_L$, $S_C$, $S_H$
    \item Compute rotation term: $R_T$
    \item Calculate final ΔE:
    \begin{equation}
        \Delta E_{00} = \sqrt{\left(\frac{\Delta L'}{S_L}\right)^2 + \left(\frac{\Delta C'}{S_C}\right)^2 + \left(\frac{\Delta H'}{S_H}\right)^2 + R_T \frac{\Delta C'}{S_C} \frac{\Delta H'}{S_H}}
    \end{equation}
\end{enumerate}

\textbf{Interpretation:}
\begin{itemize}
    \item ΔE < 1.0: Imperceptible difference (perfect match)
    \item ΔE < 2.0: Outstanding match
    \item ΔE < 3.0: Excellent match (industry standard)
    \item ΔE < 5.0: Good match (acceptable)
    \item ΔE ≥ 5.0: Poor match (noticeable difference)
\end{itemize}

\subsubsection{Recommendation Algorithm}

\textbf{Algorithm:}
\begin{enumerate}
    \item Load foundation database (187 shades)
    \item For each shade:
    \begin{itemize}
        \item Calculate ΔE using CIEDE2000
        \item Check undertone compatibility
        \item Store result with metadata
    \end{itemize}
    \item Sort results:
    \begin{itemize}
        \item Primary: Undertone-compatible shades by ΔE
        \item Secondary: Incompatible shades by ΔE
    \end{itemize}
    \item Return top 5 recommendations
    \item Include for each:
    \begin{itemize}
        \item Brand and shade name
        \item ΔE value
        \item CIELAB values
        \item Undertone
        \item MST range
        \item Match quality indicator
    \end{itemize}
\end{enumerate}

\textbf{Undertone Compatibility:}
\begin{itemize}
    \item Neutral undertone: Compatible with all shades
    \item Warm/Cool: Prefer matching undertone
    \item If insufficient matches: Include best non-matching
\end{itemize}

\section{Foundation Database}

\subsection{Database Structure}

The foundation database is stored in JSON format with the following structure:

\begin{verbatim}
[
  {
    "brand": "MAC",
    "product": "Studio Fix Fluid",
    "shade": "NC45",
    "L": 46.2,
    "a": 13.8,
    "b": 21.2,
    "undertone": "Warm",
    "mst_range": "MST-08",
    "hex": "#9B6B4F"
  },
  ...
]
\end{verbatim}

\subsection{Database Coverage}

The database includes 187 shades from 14 brands:

\begin{table}[h]
\centering
\small
\begin{tabular}{|l|c|c|}
\hline
\textbf{Brand} & \textbf{Shades} & \textbf{Percentage} \\
\hline
Fenty Beauty & 33 & 17.6\% \\
MAC & 22 & 11.8\% \\
Maybelline & 18 & 9.6\% \\
Lancôme & 14 & 7.5\% \\
NARS & 14 & 7.5\% \\
Estée Lauder & 13 & 7.0\% \\
L'Oréal & 15 & 8.0\% \\
Black Opal & 11 & 5.9\% \\
Bobbi Brown & 11 & 5.9\% \\
Pat McGrath & 10 & 5.3\% \\
Iman Cosmetics & 10 & 5.3\% \\
Hourglass & 8 & 4.3\% \\
Charlotte Tilbury & 8 & 4.3\% \\
Others & 10 & 5.3\% \\
\hline
\textbf{Total} & \textbf{187} & \textbf{100\%} \\
\hline
\end{tabular}
\caption{Foundation Database by Brand}
\end{table}

\subsection{MST Coverage}

The database provides excellent coverage across all MST classes:

\begin{table}[h]
\centering
\begin{tabular}{|c|l|c|c|}
\hline
\textbf{MST} & \textbf{Label} & \textbf{Shades} & \textbf{Percentage} \\
\hline
MST-01 & Porcelain & 23 & 12.3\% \\
MST-02 & Fair & 24 & 12.8\% \\
MST-03 & Light & 21 & 11.2\% \\
MST-04 & Light-Medium & 17 & 9.1\% \\
MST-05 & Medium & 23 & 12.3\% \\
MST-06 & Medium-Dark & 36 & 19.3\% \\
MST-07 & Tan & 31 & 16.6\% \\
MST-08 & Brown & 27 & 14.4\% \\
MST-09 & Deep Brown & 26 & 13.9\% \\
MST-10 & Deep & 20 & 10.7\% \\
\hline
\textbf{Total} & & \textbf{187} & \textbf{100\%} \\
\hline
\end{tabular}
\caption{Foundation Database by MST Class}
\end{table}

Notable features:
\begin{itemize}
    \item Excellent representation for deeper tones (MST-06 to MST-10: 140 shades, 74.9\%)
    \item Balanced coverage across all classes
    \item Multiple undertone options for each MST class
\end{itemize}

\section{Photo Quality Assessment}

\subsection{Purpose}
Photo quality significantly affects accuracy. The system includes real-time quality assessment to guide users toward better photos.

\subsection{Quality Metrics}

Six quality metrics are evaluated:

\begin{enumerate}
    \item \textbf{Brightness:}
    \begin{itemize}
        \item Measure: Mean pixel intensity
        \item Optimal: 100-180
        \item Issues: Too dark (< 80) or too bright (> 200)
    \end{itemize}
    
    \item \textbf{Blur:}
    \begin{itemize}
        \item Measure: Laplacian variance
        \item Optimal: > 100
        \item Issues: Blurry image (< 100)
    \end{itemize}
    
    \item \textbf{Color Cast:}
    \begin{itemize}
        \item Measure: RGB channel differences
        \item Optimal: Balanced channels
        \item Issues: Strong color cast (difference > 30)
    \end{itemize}
    
    \item \textbf{Face Detection:}
    \begin{itemize}
        \item Measure: Face detection success
        \item Optimal: Face detected with confidence > 0.8
        \item Issues: No face or low confidence
    \end{itemize}
    
    \item \textbf{Resolution:}
    \begin{itemize}
        \item Measure: Image dimensions
        \item Optimal: > 640×480
        \item Issues: Low resolution
    \end{itemize}
    
    \item \textbf{Exposure:}
    \begin{itemize}
        \item Measure: Histogram distribution
        \item Optimal: Well-distributed
        \item Issues: Clipping or narrow range
    \end{itemize}
\end{enumerate}

\subsection{Quality Scoring}

Overall quality score (0-100):
\begin{equation}
Score = 100 - \sum_{i=1}^{6} penalty_i
\end{equation}

where each metric contributes a penalty (0-25) based on severity.

\textbf{Quality Levels:}
\begin{itemize}
    \item Excellent: Score ≥ 90
    \item Good: Score ≥ 75
    \item Fair: Score ≥ 60
    \item Poor: Score < 60
\end{itemize}

\section{User Interfaces}

\subsection{Command-Line Interface}

The CLI provides a simple interface for batch processing and testing:

\textbf{Usage:}
\begin{verbatim}
python shadeai_cli.py image.jpg
\end{verbatim}

\textbf{Output:}
\begin{itemize}
    \item Photo quality assessment
    \item Skin tone classification (MST, ITA°)
    \item Undertone detection
    \item Top 5 shade recommendations with ΔE values
    \item Processing time
\end{itemize}

\subsection{Web Application}

The Streamlit web app provides an interactive interface:

\textbf{Features:}
\begin{itemize}
    \item File upload or camera capture
    \item Real-time quality feedback
    \item Visual results with color swatches
    \item Photo guidelines and examples
    \item Downloadable results
\end{itemize}

\textbf{Technology:}
\begin{itemize}
    \item Framework: Streamlit
    \item Backend: FastAPI
    \item Communication: HTTP requests
\end{itemize}

\subsection{REST API}

The FastAPI backend provides programmatic access:

\textbf{Endpoint:}
\begin{verbatim}
POST /api/v1/analyze
Content-Type: multipart/form-data
Body: file (image)
\end{verbatim}

\textbf{Response:}
\begin{verbatim}
{
  "skin_tone": {
    "class": "MST-08",
    "label": "Brown",
    "ita": -15.79
  },
  "undertone": {
    "label": "Neutral",
    "score": 0.1
  },
  "recommendations": [
    {
      "brand": "MAC",
      "shade": "NC45",
      "delta_e": 2.489,
      "undertone": "Warm"
    },
    ...
  ]
}
\end{verbatim}

\section{Deployment Options}

\subsection{Local Deployment}

\textbf{Requirements:}
\begin{itemize}
    \item Python 3.10+
    \item Dependencies from requirements.txt
    \item 2GB RAM minimum
\end{itemize}

\textbf{Steps:}
\begin{verbatim}
pip install -r requirements.txt
python shadeai_cli.py image.jpg
\end{verbatim}

\subsection{Docker Deployment}

\textbf{Advantages:}
\begin{itemize}
    \item Consistent environment
    \item Easy deployment
    \item Isolated dependencies
\end{itemize}

\textbf{Configuration:}
\begin{verbatim}
docker-compose up --build
# Frontend: http://localhost:8501
# Backend: http://localhost:8000
\end{verbatim}

\section{Summary}

This chapter has described the complete system design and methodology for ShadeAI, including:

\begin{itemize}
    \item Six-stage processing pipeline
    \item Detailed algorithms for each module
    \item Foundation database structure and coverage
    \item Photo quality assessment system
    \item Multiple user interfaces
    \item Deployment options
\end{itemize}

The design emphasizes modularity, accuracy, accessibility, and usability. The next chapter presents implementation details and experimental results demonstrating the system's performance.
